\input{../tex-inputs/latex-leseansicht-vorspann}

\section[Arthur Schnitzler, Richard Beer-Hofmann und Paul Goldmann an Gustav Schwarzkopf, {[}17. oder 18.{]} 8. 1900]{L04407 Arthur Schnitzler, Richard Beer-Hofmann und Paul Goldmann an Gustav Schwarzkopf, {[}17. oder 18.{]} 8. 1900}
\nopagebreak\mylabel{L04407v}
\rehead{ }\normalsize\beginnumbering\briefempfaengerindex{Schwarzkopf, Gustav@\textsc{Schwarzkopf, Gustav}!zzzGoldmann, Paul@\emph{von Paul Goldmann}!1900-08-181@{{[}17. oder 18.{]} 8. 1900}|(be}\briefempfaengerindex{Schwarzkopf, Gustav@\textsc{Schwarzkopf, Gustav}!zzzBeer-Hofmann, Richard@\emph{von Richard Beer-Hofmann}!1900-08-181@{{[}17. oder 18.{]} 8. 1900}|(be}\briefempfaengerindex{Schwarzkopf, Gustav@\textsc{Schwarzkopf, Gustav}!zzzSchnitzler, Arthur@\emph{von Arthur Schnitzler}!1900-08-181@{{[}17. oder 18.{]} 8. 1900}|(be}
\toendnotes[C]{\smallbreak\pagebreak[2]}
\correspDesc{Versand  durch Arthur Schnitzler, Richard Beer-Hofmann, Paul Goldmann im Zeitraum [17. oder 18.] 8. 1900 in Schruns
\newline{}Übermittlung  am 18. 8. 1900 in Schruns
\newline{}Zustellung  in Wien 1/1 1
\newline{}Erhalt  durch Gustav Schwarzkopf im Zeitraum [19. 8. 1900 – 23. 8. 1900?] in Tiefer Graben 23}\toendnotes[C]{\smallbreak}
\Standort{DLA, A:Schnitzler, HS.1985.1.1897.}
\physDesc{Bildpostkarte, 100 Zeichen
\newline{}Handschrift Arthur Schnitzler: Bleistift, deutsche Kurrent
\newline{}Handschrift Richard Beer-Hofmann: Bleistift, lateinische Kurrent
\newline{}Handschrift Paul Goldmann: Bleistift
\newline{}Versand: 1) Stempel: »\nobreak{}\oindex{Schruns@\textbf{Schruns}|pwk}Schruns, 18. 8. 00\nobreak{}«.   2) Stempel: »\nobreak{}\oindex{I., Innere Stadt@\textbf{I., Innere Stadt}|pwk}\textcolor{gray}{Wi}en 1/1 \textcolor{gray}{1}\nobreak{}«. }\pstart{}{\pb}Herrn Gustav Schwarzkopf\pend{}\pstart{}Wien\oindex{Wien@\textbf{Wien}|pw}\pend{}\pstart{}I. Tiefer Graben 23\oindex{Wien@\textbf{Wien}!I., Innere Stadt@\textbf{I., Innere Stadt}!Tiefer Graben 23@\textbf{Tiefer Graben 23}|pw}\pend{}{\bigskip}
\pstart
           \centering{}{\pb}\textcolor{gray}{\textbf{Schruns}}\oindex{Schruns@\textbf{Schruns}|pw}\pend
           \vspace{1em}
\pstart
           \noindent{}{\pb}Herzlich Gruß!\pend
           \pstart \spacefill\mbox{Arthur Sch}\pend{}\selectlanguage{ngerman}\vspace{1em}
\pstart
           \noindent{}{[}hs. Beer-Hofmann:{]} Herzlichst\pend
           \pstart \spacefill\mbox{Richard}\pend{}\selectlanguage{ngerman}\vspace{1em}\pstart \spacefill\mbox{{[}hs. Goldmann:{]} Dr. Goldmann}\pend{}\selectlanguage{ngerman}\endnumbering\briefempfaengerindex{Schwarzkopf, Gustav@\textsc{Schwarzkopf, Gustav}!zzzGoldmann, Paul@\emph{von Paul Goldmann}!1900-08-171@{{[}17. oder 18.{]} 8. 1900}|)be}\briefempfaengerindex{Schwarzkopf, Gustav@\textsc{Schwarzkopf, Gustav}!zzzBeer-Hofmann, Richard@\emph{von Richard Beer-Hofmann}!1900-08-171@{{[}17. oder 18.{]} 8. 1900}|)be}\briefempfaengerindex{Schwarzkopf, Gustav@\textsc{Schwarzkopf, Gustav}!zzzSchnitzler, Arthur@\emph{von Arthur Schnitzler}!1900-08-171@{{[}17. oder 18.{]} 8. 1900}|)be}\mylabel{L04407h}
\begin{anhang}
\end{anhang}\newcommand{\dateiname}{L04407}\newcommand{\titel}{Arthur Schnitzler, Richard Beer-Hofmann und Paul Goldmann an Gustav Schwarzkopf, [17. oder 18.] 8. 1900}\newcommand{\editorInnen}{Herausgegeben von Jahnke, SelmaMüller, Martin Anton}\input{../tex-inputs/latex-leseansicht-abspann}
